
\documentclass{article}
\usepackage{colortbl}
\usepackage{makecell}
\usepackage{multirow}
\usepackage{supertabular}

\begin{document}

\newcounter{utterance}

\centering \large Interaction Transcript for game `dond', experiment `coop\_it', episode 17 with gpt{-}5{-}mini{-}2025{-}08{-}07{-}no{-}reasoning{-}t1.0.
\vspace{24pt}

{ \footnotesize  \setcounter{utterance}{1}
\setlength{\tabcolsep}{0pt}
\begin{supertabular}{c@{$\;$}|p{.15\linewidth}@{}p{.15\linewidth}p{.15\linewidth}p{.15\linewidth}p{.15\linewidth}p{.15\linewidth}}
    \# & $\;$A & \multicolumn{4}{c}{Game Master} & $\;\:$B\\
    \hline

    \theutterance \stepcounter{utterance}  
    & & \multicolumn{4}{p{0.6\linewidth}}{
        \cellcolor[rgb]{0.9,0.9,0.9}{
            \makecell[{{p{\linewidth}}}]{
                \texttt{\tiny{[P1$\langle$GM]}}
                \texttt{State giocando a un gioco di negoziazione cooperativa in cui dovete accordarvi su come dividere una serie di oggetti tra voi e un altro giocatore.} \\
\\ 
\texttt{Regole:} \\
\texttt{(a) A Lei e all'altro giocatore viene dato un insieme di oggetti. Ognuno di voi riceve anche una funzione di valore segreta, che rappresenta il valore di ciascun tipo di oggetto.} \\
\texttt{(b) Si scambiano messaggi con l'altro giocatore per concordare chi si aggiudica gli oggetti. Potete inviare un massimo di 5 messaggi ciascuno, oppure terminare in anticipo facendo una proposta segreta in qualsiasi momento.} \\
\texttt{(c) A ciascuno di voi viene chiesto di inviare una proposta segreta indicando gli oggetti che desiderate, formattata tra parentesi quadre come segue: "{[}Proposta: <numero> <nome oggetto>, <numero> <nome oggetto>, <...>{]}".} \\
\texttt{(d) Se le vostre proposte sono complementari, cioè ci sono abbastanza oggetti per soddisfare entrambe le proposte, a ciascun giocatore viene assegnato un punteggio basato sulla somma dei valori degli oggetti ricevuti. In caso contrario, entrambi ricevono zero punti.} \\
\texttt{(e) L'obiettivo è massimizzare la somma del proprio punteggio e di quello dell'altro giocatore.} \\
\\ 
\texttt{Cominciamo.} \\
\\ 
\texttt{L'insieme degli oggetti disponibili è:} \\
\\ 
\texttt{1 pentola, 3 magneti, 1 piatto} \\
\\ 
\texttt{La funzione valore segreta è:} \\
\\ 
\texttt{pentola: 5, magnete: 0, piatto: 5} \\
\\ 
\texttt{IMPORTANTE: i vostri messaggi, a meno che non si tratti di una proposta segreta, vengono trasmessi direttamente all'altro giocatore, quindi non includete alcuna risposta alle regole o testo di annuncio del vostro messaggio. Per fare una proposta segreta, utilizzate il formato indicato. Non utilizzare delle parentesi quadre quando si comunica all'altro giocatore, altrimenti verrà interpretata come una proposta segreta.} \\
            }
        }
    }
    & \\ \\

    \theutterance \stepcounter{utterance}  
    & \multicolumn{4}{p{0.6\linewidth}}{
        \cellcolor[rgb]{0.9,0.9,0.9}{
            \makecell[{{p{\linewidth}}}]{
                \texttt{\tiny{[P1$\rangle$GM]}}
                \texttt{Ciao! Ecco le mie preferenze: per me pentola vale 5, piatto 5, magneti 0. Quindi non mi interessano i magneti. Se per te i magneti valgono qualcosa, possiamo fare uno scambio efficiente.} \\
\\ 
\texttt{Propongo: io prendo pentola e piatto, tu prendi i 3 magneti. Se per te i magneti valgono poco e preferisci uno tra pentola/piatto, dimmelo e vediamo se c’è un’alternativa che massimizza la somma.} \\
            }
        }
    }
    & & \\ \\

    \theutterance \stepcounter{utterance}  
    & & & \multicolumn{2}{p{0.3\linewidth}}{
        \cellcolor[rgb]{0.9,0.9,0.9}{
            \makecell[{{p{\linewidth}}}]{
                \texttt{\tiny{[GM$|$GM]}}
                \texttt{continue} \\
            }
        }
    }
    & & \\ \\

    \theutterance \stepcounter{utterance}  
    & & & \multicolumn{2}{p{0.3\linewidth}}{
        \cellcolor[rgb]{0.9,0.9,0.9}{
            \makecell[{{p{\linewidth}}}]{
                \texttt{\tiny{[GM$|$GM]}}
                \texttt{Ciao! Ecco le mie preferenze: per me pentola vale 5, piatto 5, magneti 0. Quindi non mi interessano i magneti. Se per te i magneti valgono qualcosa, possiamo fare uno scambio efficiente.} \\
\\ 
\texttt{Propongo: io prendo pentola e piatto, tu prendi i 3 magneti. Se per te i magneti valgono poco e preferisci uno tra pentola/piatto, dimmelo e vediamo se c’è un’alternativa che massimizza la somma.} \\
            }
        }
    }
    & & \\ \\

    \theutterance \stepcounter{utterance}  
    & & \multicolumn{4}{p{0.6\linewidth}}{
        \cellcolor[rgb]{0.9,0.9,0.9}{
            \makecell[{{p{\linewidth}}}]{
                \texttt{\tiny{[GM$\rangle$P2]}}
                \texttt{State giocando a un gioco di negoziazione cooperativa in cui dovete accordarvi su come dividere una serie di oggetti tra voi e un altro giocatore.} \\
\\ 
\texttt{Regole:} \\
\texttt{(a) A Lei e all'altro giocatore viene dato un insieme di oggetti. Ognuno di voi riceve anche una funzione di valore segreta, che rappresenta il valore di ciascun tipo di oggetto.} \\
\texttt{(b) Si scambiano messaggi con l'altro giocatore per concordare chi si aggiudica gli oggetti. Potete inviare un massimo di 5 messaggi ciascuno, oppure terminare in anticipo facendo una proposta segreta in qualsiasi momento.} \\
\texttt{(c) A ciascuno di voi viene chiesto di inviare una proposta segreta indicando gli oggetti che desiderate, formattata tra parentesi quadre come segue: "{[}Proposta: <numero> <nome oggetto>, <numero> <nome oggetto>, <...>{]}".} \\
\texttt{(d) Se le vostre proposte sono complementari, cioè ci sono abbastanza oggetti per soddisfare entrambe le proposte, a ciascun giocatore viene assegnato un punteggio basato sulla somma dei valori degli oggetti ricevuti. In caso contrario, entrambi ricevono zero punti.} \\
\texttt{(e) L'obiettivo è massimizzare la somma del proprio punteggio e di quello dell'altro giocatore.} \\
\\ 
\texttt{Cominciamo.} \\
\\ 
\texttt{L'insieme degli oggetti disponibili è:} \\
\\ 
\texttt{1 pentola, 3 magneti, 1 piatto} \\
\\ 
\texttt{La funzione valore segreta è:} \\
\\ 
\texttt{pentola: 4, magnete: 2, piatto: 0} \\
\\ 
\texttt{IMPORTANTE: i vostri messaggi, a meno che non si tratti di una proposta segreta, vengono trasmessi direttamente all'altro giocatore, quindi non includete alcuna risposta alle regole o testo di annuncio del vostro messaggio. Per fare una proposta segreta, utilizzate il formato indicato. Non utilizzare delle parentesi quadre quando si comunica all'altro giocatore, altrimenti verrà interpretata come una proposta segreta.} \\
\\ 
\\ 
\texttt{Ciao! Ecco le mie preferenze: per me pentola vale 5, piatto 5, magneti 0. Quindi non mi interessano i magneti. Se per te i magneti valgono qualcosa, possiamo fare uno scambio efficiente.} \\
\\ 
\texttt{Propongo: io prendo pentola e piatto, tu prendi i 3 magneti. Se per te i magneti valgono poco e preferisci uno tra pentola/piatto, dimmelo e vediamo se c’è un’alternativa che massimizza la somma.} \\
            }
        }
    }
    & \\ \\

    \theutterance \stepcounter{utterance}  
    & & & \multicolumn{4}{p{0.6\linewidth}}{
        \cellcolor[rgb]{0.9,0.9,0.9}{
            \makecell[{{p{\linewidth}}}]{
                \texttt{\tiny{[GM$\langle$P2]}}
                \texttt{Grazie per la trasparenza, siamo allineati sull’obiettivo.} \\
\\ 
\texttt{Le mie valutazioni: pentola 4, magnete 2 ciascuno, piatto 0. Quindi:} \\
\texttt{{-} A me valgono i magneti (tot 6) e anche la pentola (4), il piatto per me vale 0.} \\
\texttt{{-} A te valgono pentola (5) e piatto (5), i magneti 0.} \\
\\ 
\texttt{La somma massima possibile è 5 (pentola) + 5 (piatto) + 3×2 (magneti) = 16, che otteniamo se:} \\
\texttt{{-} Tu prendi pentola e piatto} \\
\texttt{{-} Io prendo i 3 magneti} \\
\\ 
\texttt{Questo coincide con la tua proposta ed è socialmente efficiente. Se sei d’accordo, confermo così. Se preferisci una piccola variazione per bilanciare la percezione (anche se il punteggio totale resta 16), posso rinunciare alla pentola e tenermi solo i magneti.} \\
\\ 
\texttt{Se sei d’accordo a chiudere adesso, procedo con la proposta segreta complementare.} \\
            }
        }
    }
     \\ \\

    \theutterance \stepcounter{utterance}  
    & & & \multicolumn{2}{p{0.3\linewidth}}{
        \cellcolor[rgb]{0.9,0.9,0.9}{
            \makecell[{{p{\linewidth}}}]{
                \texttt{\tiny{[GM$|$GM]}}
                \texttt{continue} \\
            }
        }
    }
    & & \\ \\

    \theutterance \stepcounter{utterance}  
    & & & \multicolumn{2}{p{0.3\linewidth}}{
        \cellcolor[rgb]{0.9,0.9,0.9}{
            \makecell[{{p{\linewidth}}}]{
                \texttt{\tiny{[GM$|$GM]}}
                \texttt{Grazie per la trasparenza, siamo allineati sull’obiettivo.} \\
\\ 
\texttt{Le mie valutazioni: pentola 4, magnete 2 ciascuno, piatto 0. Quindi:} \\
\texttt{{-} A me valgono i magneti (tot 6) e anche la pentola (4), il piatto per me vale 0.} \\
\texttt{{-} A te valgono pentola (5) e piatto (5), i magneti 0.} \\
\\ 
\texttt{La somma massima possibile è 5 (pentola) + 5 (piatto) + 3×2 (magneti) = 16, che otteniamo se:} \\
\texttt{{-} Tu prendi pentola e piatto} \\
\texttt{{-} Io prendo i 3 magneti} \\
\\ 
\texttt{Questo coincide con la tua proposta ed è socialmente efficiente. Se sei d’accordo, confermo così. Se preferisci una piccola variazione per bilanciare la percezione (anche se il punteggio totale resta 16), posso rinunciare alla pentola e tenermi solo i magneti.} \\
\\ 
\texttt{Se sei d’accordo a chiudere adesso, procedo con la proposta segreta complementare.} \\
            }
        }
    }
    & & \\ \\

    \theutterance \stepcounter{utterance}  
    & & \multicolumn{4}{p{0.6\linewidth}}{
        \cellcolor[rgb]{0.9,0.9,0.9}{
            \makecell[{{p{\linewidth}}}]{
                \texttt{\tiny{[P1$\langle$GM]}}
                \texttt{Grazie per la trasparenza, siamo allineati sull’obiettivo.} \\
\\ 
\texttt{Le mie valutazioni: pentola 4, magnete 2 ciascuno, piatto 0. Quindi:} \\
\texttt{{-} A me valgono i magneti (tot 6) e anche la pentola (4), il piatto per me vale 0.} \\
\texttt{{-} A te valgono pentola (5) e piatto (5), i magneti 0.} \\
\\ 
\texttt{La somma massima possibile è 5 (pentola) + 5 (piatto) + 3×2 (magneti) = 16, che otteniamo se:} \\
\texttt{{-} Tu prendi pentola e piatto} \\
\texttt{{-} Io prendo i 3 magneti} \\
\\ 
\texttt{Questo coincide con la tua proposta ed è socialmente efficiente. Se sei d’accordo, confermo così. Se preferisci una piccola variazione per bilanciare la percezione (anche se il punteggio totale resta 16), posso rinunciare alla pentola e tenermi solo i magneti.} \\
\\ 
\texttt{Se sei d’accordo a chiudere adesso, procedo con la proposta segreta complementare.} \\
            }
        }
    }
    & \\ \\

    \theutterance \stepcounter{utterance}  
    & \multicolumn{4}{p{0.6\linewidth}}{
        \cellcolor[rgb]{0.9,0.9,0.9}{
            \makecell[{{p{\linewidth}}}]{
                \texttt{\tiny{[P1$\rangle$GM]}}
                \texttt{Perfetto, concordo con la divisione efficiente: tu prendi i 3 magneti, io pentola e piatto. Procedo con la proposta segreta complementare.} \\
\\ 
\texttt{{[}Proposta: 1 pentola, 1 piatto{]}} \\
            }
        }
    }
    & & \\ \\

    \theutterance \stepcounter{utterance}  
    & & & \multicolumn{2}{p{0.3\linewidth}}{
        \cellcolor[rgb]{0.9,0.9,0.9}{
            \makecell[{{p{\linewidth}}}]{
                \texttt{\tiny{[GM$|$GM]}}
                \texttt{proposal} \\
            }
        }
    }
    & & \\ \\

    \theutterance \stepcounter{utterance}  
    & & & \multicolumn{2}{p{0.3\linewidth}}{
        \cellcolor[rgb]{0.9,0.9,0.9}{
            \makecell[{{p{\linewidth}}}]{
                \texttt{\tiny{[GM$|$GM]}}
                [1, 0, 1]
            }
        }
    }
    & & \\ \\

    \theutterance \stepcounter{utterance}  
    & & \multicolumn{4}{p{0.6\linewidth}}{
        \cellcolor[rgb]{0.9,0.9,0.9}{
            \makecell[{{p{\linewidth}}}]{
                \texttt{\tiny{[GM$\rangle$P2]}}
                \texttt{L'altro giocatore ha presentato una proposta segreta. Ora tocca a Lei presentare una sua proposta.} \\
            }
        }
    }
    & \\ \\

    \theutterance \stepcounter{utterance}  
    & & & \multicolumn{4}{p{0.6\linewidth}}{
        \cellcolor[rgb]{0.9,0.9,0.9}{
            \makecell[{{p{\linewidth}}}]{
                \texttt{\tiny{[GM$\langle$P2]}}
                \texttt{{[}Proposta: 1 pentola, 0 magneti, 1 piatto{]}} \\
            }
        }
    }
     \\ \\

    \theutterance \stepcounter{utterance}  
    & & & \multicolumn{2}{p{0.3\linewidth}}{
        \cellcolor[rgb]{0.9,0.9,0.9}{
            \makecell[{{p{\linewidth}}}]{
                \texttt{\tiny{[GM$|$GM]}}
                \texttt{proposal} \\
            }
        }
    }
    & & \\ \\

    \theutterance \stepcounter{utterance}  
    & & & \multicolumn{2}{p{0.3\linewidth}}{
        \cellcolor[rgb]{0.9,0.9,0.9}{
            \makecell[{{p{\linewidth}}}]{
                \texttt{\tiny{[GM$|$GM]}}
                [1, 0, 1]
            }
        }
    }
    & & \\ \\

    \theutterance \stepcounter{utterance}  
    & & & \multicolumn{2}{p{0.3\linewidth}}{
        \cellcolor[rgb]{0.9,0.9,0.9}{
            \makecell[{{p{\linewidth}}}]{
                \texttt{\tiny{[GM$|$GM]}}
                [[1, 0, 1], [1, 0, 1]]
            }
        }
    }
    & & \\ \\

\end{supertabular}
}

\end{document}
